\documentclass[10pt]{article}
\usepackage[margin=20mm]{geometry}
\usepackage{amsmath,amssymb,mathtools,fontspec,unicode-math,xurl,hyperref}
\IfFontExistsTF{DejaVu Serif}{\setmainfont{DejaVu Serif}}{\setmainfont{Latin Modern Roman}}
\setmathfont{latinmodern-math.otf}
\setlength{\parindent}{0pt}\setlength{\parskip}{6pt}\sloppy
\title{Technical verification of a regular unimodular
triangulation of matroid base polytopes}
\author{AI-assisted technical review}
\date{}
\begin{document}
\maketitle
Abstract: The deletion–contraction construction is supported by direct checks of the resonance-flat argument, primitive sublattices, finite perturbations and the unit-height joining vector. The extension to integral generalized permutahedra admits an explicit matroid-rank description on each unit cube and a common regular height across the cubes. An exact check also verifies all forty-two maximal simplices in the stated \(K_4\) example. One auxiliary concern is confirmed: the alcoved-polytope comparison in Remark 3.4 uses an incompatible ambient lattice. Explicit integral triangles have a nonintegral intersection, and primitive face segments give an empty Cayley tetrahedron of normalized volume three. This does not enter the main theorem or its two corollaries. The assessment is Minor, with scientific adjudication pending.\par
\section*{Identity and reviewed scope}
The reviewed source is S. Backman and G. Liu, A Regular Unimodular Triangulation of the Matroid Base Polytope, arXiv:2309.10229v2, 4 April 2024, twelve pages [paper]. The frozen journal identity is Duke Mathematical Journal 174 (2025), no. 14, 3047–3061, DOI https://doi.org/10.1215/00127094-2025-0015. The acquired author preprint has twelve pages; the fifteen-page journal version was not acquired. Equality of their mathematical contents is not assumed. All page numbers below refer to the acquired PDF, whose printed numbering agrees with its PDF numbering.\par
The review covers the definitions, Theorems 3.1 and 3.5, Corollaries 3.6 and 3.7, Remark 3.4 and Example 3.8. It distinguishes a unimodular triangulation from a flag one; the article does not claim to settle White’s conjecture. Unimodularity is measured in \(\operatorname{lin}(P)\cap\mathbb Z^n\), as in Definition 2.10. This relative lattice is essential for lower-dimensional polytopes and for the concern below.\par
The full PDF text, complete released TeX and bibliography were read. All twelve primary-source pages were actually inspected. The fourteen manifest inputs, totaling 1,239,502 bytes, matched their current recorded sizes and SHA-256 hashes. Three native bounded checks supplied lattice, genericity and extension arguments, with roles and actual page inspections separately recorded. The reviewed author TeX was read as evidence and was not executed.\par
\section*{Verification of the principal construction}
\subsection*{Affine spans and saturated direction lattices}
Lemma 2.6, p. 4, can be checked without a hidden full-dimensional assumption. Form a graph on \([n]\) whose edges are the pairs appearing in the edge directions of the polytope. If its connected components are \(C_1,\ldots,C_j\), its edge vectors span exactly \[V=\{z\in\mathbb R^n:z_{C_i}=0\text{ for every }i\}.\] For an integral vertex \(v\), the affine span is \(v+V\). Taking \(S_i=C_1\cup\cdots\cup C_i\) gives the asserted flag equations, with integral right-hand sides \(v_{S_i}\). Isolated graph vertices impose fixed-coordinate equations, so points, loops and coloops present no exception to this argument.\par
For each component, choose a spanning tree and a root. Every difference \(e_a-e_b\) in that component is an integral sum of oriented tree edges, and every integral vector with component sum zero is an integral sum of such differences. Thus any collection of type-\(A\) root directions generates the entire lattice in its real span. This gives Lemma 3.3, p. 6, including dependent pairs of direction spaces.\par
In Lemma 3.2 the intended independence condition is \(V\cap W=\{0\}\). For a rational subspace \(X\subset V\), the subgroup \(X\cap\mathbb Z^n\) is primitive in \(V\cap\mathbb Z^n\): if an integral multiple of an integral vector belongs to \(X\), the vector itself belongs to the real subspace \(X\). Its lattice basis can therefore be extended to one for \(V\cap\mathbb Z^n\). The same applies to \(Y\subset W\). Complementarity and independence make the two extended bases a basis of \((V+W)\cap\mathbb Z^n\). Uniqueness of its coordinates then proves that the two smaller bases generate \((X+Y)\cap\mathbb Z^n\). The source’s empty-set symbol in its independence definition is an evident typesetting error; the proof uses zero intersection.\par
\subsection*{Genericity, refinement and the joining lattice}
Let the two coordinate slices be \(P_0\) and \(P_1\). The initial height \(g\) is affine on each slice. A lower supporting affine function therefore cuts out a face \(F_i\) on each side, and its cell is their Cayley join. If \[L=\operatorname{lin}(F_0)\cap\operatorname{lin}(F_1)\] had positive dimension, Lemma 2.6 would make it a resonance flat. The linear part of the same supporting affine function restricts to both slice heights on \(L\). Consequently \((\ell_0-\ell_1)|_L\) would be constant, contrary to Definition 2.12. Hence \(L=\{0\}\). This verifies the decisive genericity step on p. 7.\par
The small-parameter choice can be made uniform. There are only finitely many candidate affine bases and support inequalities in the finite vertex configuration. For a height \(g+\epsilon h\), every relevant residual is \(\alpha+\epsilon\beta\). All nonzero \(\alpha\) retain their signs for one sufficiently small positive \(\epsilon\); zero residuals are resolved by \(h\). Equivalently, the perturbed height lies in a secondary cone whose closure contains \(g\). Its subdivision refines that of \(g\). The restrictions to the two coordinate faces are exactly the prescribed regular triangulations, because adding an affine function and multiplying by a positive scalar preserve them.\par
A maximal refined cell is thus \[T=\operatorname{conv}\bigl((T_0,0)\cup(T_1,1)\bigr),\] where each \(T_i\) is a unimodular simplex in \(F_i\). Independence of the two direction spaces makes this a simplex. Lemmas 3.2 and 3.3 provide a basis \(B\) for the integral horizontal direction lattice using edges from \(T_0\) and \(T_1\). Take any joining vector \(w\) from level zero to level one. For \(z\in\operatorname{lin}(T)\cap\mathbb Z^n\), subtracting \(z_nw\) leaves an integral horizontal vector, hence an integral combination of \(B\). Thus \(B\cup\{w\}\) is a full lattice basis, proving unimodularity. If one side is empty, induction applies to the other side. A singleton, including the empty-ground-set matroid if it is allowed, has its unique zero-dimensional unimodular triangulation.\par
Theorem 3.5 follows by repeating these finite perturbations. For example, after an irrelevant common scaling, set \(\epsilon_k=t^k\). Each finite support residual is a polynomial in \(t\); its first nonzero coefficient fixes its sign for small positive \(t\). This supplies one compatible hierarchy for the finitely many branches.\par
\subsection*{The unit-cube extension}
Here is an explicit verification of the step in Corollary 3.6, p. 9. Write an integral generalized permutahedron as a base polytope for an integral submodular function \(f\), normalized at the empty set. In a nonempty unit-cube cell, put \(x=a+y\) with \(a\in\mathbb Z^n\) and \(0\le y_i\le1\), and set \(h(S)=f(S)-a_S\). Feasibility gives \(h(S)\ge0\). Define \[r(S)=\min_{T\subseteq[n]}\bigl(h(T)+|S\setminus T|\bigr).\] The function \(r\) is integral, normalized, monotone and submodular, and \(r(S)\le |S|\). For submodularity, take minimizing sets for two arguments and use submodularity of \(h\), together with the corresponding elementary inequality for the two set differences. For a feasible \(y\), the inequalities \(h(T)+|[n]\setminus T|\ge y_{[n]}=h([n])\) show that \(r([n])=h([n])\). These are the rank axioms of a matroid.\par
Moreover, the original cube cell is exactly \(B(r)\). The forward inclusion follows from \(y_S\le y_T+|S\setminus T|\le h(T)+|S\setminus T|\). Conversely, \(r(S)\le h(S)\) gives the original submodular inequalities, while monotonicity and \(r(\{i\})\le1\) give \(0\le y_i\le1\) for a vector of total sum \(r([n])\). This proves the required translated-matroid description.\par
One global height realizing the dicing is \[g(z)=\sum_{i=1}^n z_i^2,\qquad z\in P\cap\mathbb Z^n.\] On \([a_i,a_i+1]\), the linear interpolation of \(z_i^2\) has slope \(2a_i+1\); the sum is convex and affine on each cube. The preceding integrality argument supplies its cell vertices. The common perturbation hierarchy in Corollary 3.6 therefore refines this regular subdivision and triangulates every cell compatibly. Only finitely many cells occur, so one set of small parameters works for all of them.\par
For Corollary 3.7, the map \[(x_1,\ldots,x_n)\longmapsto
 \bigl(r(E)-\textstyle\sum_i x_i,x_1,\ldots,x_n\bigr)\] has an integral inverse on its affine image. Its linear part identifies the relative lattices. The image of the independence polytope has only edge directions \(e_i-e_j\), including the directions \(e_i-e_0\) from adding an independent element, and has integral vertices. It is an integral generalized permutahedron. The map preserves regularity and unimodularity of the constructed triangulation.\par
\section*{A confirmed auxiliary concern}
Problem 1 (The lattice in Remark 3.4: Minor).  On p. 8, an alcoved polytope is placed in \(H=\{x\in\mathbb R^n:\sum_i x_i=0\}\), while the surrounding definitions continue to use the ambient integer lattice. The remark claims integrality of intersections and unimodular triangulability of the two-polytope Cayley construction in that setting. These assertions fail with the stated lattice.\par
The facet directions in the remark are interpreted as the usual type-\(A\) facet normals; this is also the interpretation required by its subsequent assertion about \(\operatorname{lin}(P)^\perp\). Work in \(n=3\), and use the coordinates \[u=x_1-x_2,\qquad v=x_2-x_3.\] Their inverse on \(H\) is \[x=\left(\frac{2u+v}{3},\frac{v-u}{3},
             -\frac{u+2v}{3}\right).\] Thus the lattice \(\Lambda=H\cap\mathbb Z^3\) corresponds to \[\{(u,v)\in\mathbb Z^2:u+2v\equiv0\pmod3\},\] not to all of \(\mathbb Z^2\).\par
Let \[P=\{u\ge0,\ v\ge0,\ u+v\le3\}
   =\operatorname{conv}\{0,a,b\},
 \quad a=(2,-1,-1),\quad b=(1,1,-2),\] and let \(Q=P+(1,-1,0)\). Both are integral triangles in \(H\). Their three facet normals are roots, up to signs: \(e_1-e_2\), \(e_2-e_3\), and \(e_1-e_3\). The translation changes \((u,v)\) by \((2,-1)\). Consequently \[P\cap Q=\{u\ge2,\ v\ge0,\ u+v\le3\}.\] Its vertex \((u,v)=(2,0)\) is \[(4/3,-2/3,-2/3),\] which is not integral. This is a bounded, explicit counterexample to statement (B) in the remark’s asserted alcoved setting.\par
The Cayley claim has an equally concrete obstruction on the face-closed family. Take the primitive face segments \(S_a=[0,a]\) and \(S_b=[0,b]\). Each is a unimodular one-dimensional simplex. In the lattice coordinates \((x_1,x_2)\) for \(\Lambda\), their direction vectors are \((2,-1)\) and \((1,1)\). The Cayley tetrahedron \[C=\operatorname{conv}\bigl((S_a,0)\cup(S_b,1)\bigr)\] has normalized volume \[\left|\det\begin{pmatrix}
 2&0&1\\[-1mm]-1&0&1\\0&1&1
 \end{pmatrix}\right|=3.\] Its only lattice points are its four vertices: an integral height in \([0,1]\) is zero or one, and both face segments are primitive. A triangulation using lattice vertices cannot subdivide this empty tetrahedron further, so it cannot be unimodular. This also disproves the relevant two-factor version of (A) with the literal ambient lattice.\par
There is substantial counterevidence to a broader conclusion. The root-direction saturation in Lemma 3.3 is correct; it concerns edge directions of matroid polytopes. Theorem 3.1 and its extensions do not use Remark 3.4. The obstruction instead arises from carrying a statement about root normals into a different lattice realization.\par
A repair must specify the lattice as well as the real hyperplane. For example, the usual difference-coordinate realization uses the weight lattice \[\Lambda_{\mathrm w}
 =\{x\in H:x_i-x_j\in\mathbb Z\},\] which maps bijectively to \(\mathbb Z^2\) in the displayed coordinates. It contains \(\Lambda\) with index three. The fractional intersection vertex is a lattice point in \(\Lambda_{\mathrm w}\), and the displayed segments acquire additional lattice points there. Alternatively one may work in the usual coordinate realization with its stipulated integer lattice. These changes remove the counterexamples to the intended alcoved comparison, but they change the lattice used by the printed formulation. They cannot be suppressed as an unimodular coordinate change of \(\Lambda\). The ambient statement that \(\operatorname{lin}(P)^\perp\) is spanned by roots also needs its common diagonal direction or quotient convention specified: for a full-dimensional \(P\subset H\), \(\operatorname{lin}(P)^\perp=\mathbb R(1,1,1)\) in \(\mathbb R^3\).\par
The acquired primary reference confirms the convention behind this repair: Haase–Paffenholz–Piechnik–Santos, arXiv:1405.1687v3, p. 36, defines a type-\(\Gamma\) polytope using the lattice dual to the root lattice; Lemma 4.15 on pp. 72–73 then concerns that lattice and the canonical triangulations [HPPS]. Those pages were actually inspected. The acquired reference is an eighty-nine-page preprint, not an asserted copy of the 2021 published monograph. This comparison supports the lattice correction without treating the cited lemma as false.\par
The established impact is on the auxiliary comparison and its claimed alternative proof, not the article’s new triangulation theorem. That dependency distinction is the reason for the Minor assessment; the length of a possible repair is not the severity criterion.\par
\section*{Exact example checks and clarifications}
The sixteen basis triples in Example 3.8, p. 10, are exactly the spanning trees of \(K_4\) after correcting the repeated edge label \(v_0v_1=2\) to \(v_0v_2=2\). The four missing triples are \(\{0,1,2\}\), \(\{0,3,4\}\), \(\{1,3,5\}\) and \(\{2,4,5\}\), the four triangle circuits.\par
The forty-two listed six-vertex cells were checked with exact integer and rational arithmetic. Projecting away the last coordinate identifies the affine lattice \(\sum x_i=3\) with \(\mathbb Z^5\). For every listed cell, the determinant of its affine vertex matrix has absolute value one. Solving for its supporting height in each of the five \(\epsilon\) coefficients gives a strictly positive lexicographic residual at every other vertex. Thus each listed simplex is an actual lower facet for one sufficiently small common hierarchy. This computation uses the printed positive final coefficient in the functional; the all-negative interpretation also yields the same forty-two facets.\par
As an independent completeness check, the lattice-point counts for \(tP\), for \(t=0,\ldots,5\), are \[1,\ 16,\ 101,\ 396,\ 1167,\ 2848.\] They were enumerated using \(0\le x_i\le t\), \(\sum x_i=3t\), and the four triangle inequalities \(x_C\le2t\). The fifth finite difference is \(42\), hence the normalized volume is \(42\). The forty-two distinct unimodular lower simplices therefore exhaust the polytope’s volume; no further full-dimensional lower cell remains.\par
Genericity of the specified functionals was also checked in dimensions one through five. All independent collections of \(m-1\) nonzero zero–one normal vectors were tested using exact cofactor kernels. The numbers of independent collections were \(1,3,21,430,27990\); none had zero pairing with the printed coefficient vector. Every positive-dimensional flat contains a line flat, so this verifies the needed genericity in this example. It does not assert that one fixed exponential prescription works in every dimension.\par
This particular triangulation is not flag. For instance, the vertex indices \(\{0,6,8\}\) form a clique: each pair is an edge of a listed simplex, but no listed simplex contains the triple. This is a check of the exhibited triangulation only. It does not prove that all choices in Theorem 3.5 fail to be flag.\par
The introductory volume discussion on p. 1 needs consistent normalization. If \(d=\dim P\), the number of maximal unimodular simplices and \(h^*(1)\) equal \(\operatorname{Vol}_{\mathbb Z}(P)\), whereas \[[t^d]\,L_P(t)=\operatorname{Vol}_{\mathbb Z}(P)/d!.\] For \(P_{U_{1,3}}=\operatorname{conv}(e_1,e_2,e_3)\) there is one unimodular simplex and \(h^*=1\), but \(L_P(t)=\binom{t+2}{2}\) has leading coefficient \(1/2\). This is recorded as a normalization clarification. It changes none of the triangulation arguments. The misplaced \(f_1\) subscript on the \(P_0\) restriction on p. 7 likewise should be \(f_0\), as the defining height already says. These visible local corrections are not additional substantive findings.\par
\section*{Assessment and state}
The main proof and both extensions are supported within the reviewed version. Standard matroid base-polytope characterizations and finite regular-subdivision theory are identified dependencies; no complete formal verification of all cited literature is claimed. For the general face- and translation-closed families in Remark 3.4, the forward implications from lattice complementarity and the two-subspace duality have direct bounded proofs. The two arbitrary-family converse implications were not independently reconstructed in this review. This is a verification limit, not a claimed counterexample to those equivalences. The accessible search confirmed the pinned arXiv version. No separate correction, supplement or code release resolving the lattice issue was identified. The conference version and journal version are distinct bibliographic objects; no equality is inferred from their titles. Absence of code is not a mathematical defect.\par
Item | Assessment | Verification and remaining work\par
Remark 3.4 lattice | Minor | Explicit intersection and Cayley counterexamples; specify the intended lattice. Author adoption unknown.\par
Main construction | Supported | Genericity, saturation, relative lattice and common perturbations checked.\par
Extensions | Supported | Explicit cube-cell rank and integral independence-polytope lift supplied.\par
Example 3.8 | Supported | All forty-two cells and their completeness checked exactly; this triangulation is not flag.\par
Other displays | Clarification | Volume convention and visible subscripts/labels; no extra severity count.\par
Scientific decision remains pending, with no independent acceptance or human-review attestation.\par
\section*{Supporting checks}
ca9be314b2387e5d3c8ceb4397996efe426493f238d8cad834271b85e3abc492.\par
The exact computation includes all affine support residuals, cell determinants, lattice-point counts and genericity checks. A failed first parsing attempt is preserved and contributes no scientific result; its corrected successor was actually run.\par
Three existing native helpers performed bounded checks only. Actual viewing receipts are distinct from rendering records; the latter retain their original false inspection flags. The delivery state retains `scientific\_decision=pending` and `human\_reviews=[]`.\par
9 S. Backman and G. Liu, A Regular Unimodular Triangulation of the Matroid Base Polytope, arXiv:2309.10229v2, 4 April 2024, https://arxiv.org/abs/2309.10229v2. C. Haase, A. Paffenholz, L. Piechnik and F. Santos, Existence of unimodular triangulations—positive results, arXiv:1405.1687v3. Acquired primary preprint; lattice definition and Lemma 4.15 inspected in bounded ranges, https://arxiv.org/abs/1405.1687v3.\par
\end{document}
